\documentclass[12pt]{article}
\usepackage[utf8]{inputenc}
\usepackage[T1]{fontenc}
\usepackage{hyperref}
\title{Rundi Language Resources for K-12 Curriculum Development}
\author{Compiled by Claude Sonnet 5.5 (claude-sonnet-5-5) for Andr\'as Kornai}
\date{}
\begin{document}
\maketitle

\section{Introduction}
Rundi (Kirundi; ISO 639-3 run) is the national language of Burundi and is used in schooling there. My search found mostly learner materials aimed at foreigners. It did not find official Burundian curriculum documents or government textbooks, which should be sought from the Ministry of National Education and local publishers.

\section{Primary Educational Resources}
\begin{itemize}
\item Kirundi Children's Dictionary, an illustrated Kirundi-English, English-Kirundi dictionary of about 150 words (2014, self-published): \url{https://www.langcen.cam.ac.uk/resources/langr/data/rd_102.html}
\item Kirundi Basic Course, student text (a reissue of a 1960s FSI course), listed by the Cambridge Language Centre: \url{https://www.langcen.cam.ac.uk/resources/langr/rundi.html}
\end{itemize}

\section{Linguistic and Reference Materials}
\begin{itemize}
\item Kirundi-English/English-Kirundi Dictionary \&amp; Phrasebook (Callixte Nizigama, Hippocrene, 2020; about 4,000 entries): \url{https://www.semcoop.com/kirundi-english-english-kirundi-dictionary-phrasebook}
\item Online Kirundi dictionary entries (matana.de): \url{https://www.matana.de/index1.php?l=3&amp;q=rain}
\item Further Cambridge Language Centre listing: \url{https://www.langcen.cam.ac.uk/resources/langr/data/rd_103.html}
\end{itemize}

\section{Religious and Community Resources}
None found in this search.

\section{Digital and Audio Resources}
The Cambridge Language Centre notes that audio is available with the basic course (see the link above). No other digital resources were found.

\section{Curriculum Development Context}
Kirundi is a language of instruction in Burundi, so a curriculum can build on existing national school materials. Those were not found in search and need a direct look at Ministry of Education and Burundian publisher sources. The resources above supply vocabulary and learner-level support.
\end{document}